\documentclass[12pt,letterpaper]{article}
\usepackage[utf8]{inputenc}
\usepackage[spanish]{babel}
\usepackage{amsmath}
\usepackage{amsfonts}
\usepackage{amssymb}
\usepackage[margin=2.5cm]{geometry}

\begin{document}

\begin{center}
    \Large \textbf{Evaluación de Examen 2: Unidad II} \\
    \large Física para Ciencias de la Salud (005-1713) \\
    \vspace{0.5cm}
    \textbf{Estudiante:} Deliangelis Inojosa \quad \textbf{C.I.:} 31.516.750
    \vspace{0.3cm} \\
    \large \textbf{Calificación Final: 9,5/10 puntos}
\end{center}

\vspace{0.5cm}

\section*{Análisis de Resultados}

\subsection*{1. Cinemática (Aceleración media)}
\textbf{Procedimiento y resultado:} Extrae los datos a la perfección y plantea la ecuación de aceleración de forma correcta. Sustituye y calcula $0,6 \text{ m/s} / 0,15 \text{ s}$, llegando al resultado exacto de $4 \text{ m/s}^2$. La redacción final de su respuesta es muy completa. \\
\textbf{Veredicto:} \textbf{Correcto} (2/2 pts).

\subsection*{2. Dinámica (Leyes de Newton)}
\textbf{Procedimiento y resultado:} Muestra un excelente manejo del álgebra al plantear la Segunda Ley de Newton y despejar analíticamente la aceleración ($a = \frac{F}{m}$). Aunque no plasmó explícitamente los números en la fracción durante el desarrollo, su conclusión final de $6 \text{ m/s}^2$ demuestra que el cálculo fue efectuado correctamente. \\
\textbf{Veredicto:} \textbf{Correcto} (2/2 pts).

\subsection*{3. Trabajo Mecánico}
\textbf{Procedimiento y resultado:} Justifica de gran manera el uso del ángulo ($0^\circ$) al indicar que tanto el movimiento como la fuerza van hacia arriba, aplicando $\cos(0^\circ) = 1$. La multiplicación ($400 \text{ N} \cdot 0,8 \text{ m}$) es perfecta y su resultado de $320 \text{ Joules}$ es completamente exacto. \\
\textbf{Veredicto:} \textbf{Correcto} (2/2 pts).

\subsection*{4. Energía Potencial}
\textbf{Procedimiento y resultado:} Extrae los datos (masa, altura y gravedad) y procede al cálculo en su apartado de "Desarrollo", multiplicando los tres factores ($1,2 \text{ kg} \cdot 9,8 \text{ m/s}^2 \cdot 1,5 \text{ m}$). El valor obtenido de $17,64 \text{ Joules}$ es verificado y correcto. \\
\textbf{Veredicto:} \textbf{Correcto} (2/2 pts).

\subsection*{5. Potencia Mecánica}
\textbf{Procedimiento y resultado:} La estudiante extrae los datos correspondientes. No obstante, en la zona de desarrollo y fórmula omite por completo indicar qué ecuación física rige el fenómeno o mostrar la división correspondiente ($75 / 60$). Salta directamente a declarar la respuesta ($1,25 \text{ Watts}$); aunque el número es el correcto, un examen de desarrollo requiere que se evidencie de dónde provienen los cálculos. \\
\textbf{Veredicto:} \textbf{Parcialmente correcto} (Se penaliza de forma leve la omisión del procedimiento matemático/fórmula) (1,5/2 pts).

\vspace{0.5cm}
\section*{Comentarios Generales y Sugerencias}
¡Felicidades por un examen sobresaliente! La estudiante cuenta con una presentación pulcra y una estrategia de resolución sumamente clara dividiendo cada página en \textit{Datos}, \textit{Fórmula}, \textit{Desarrollo} y \textit{Resultado}.
\begin{itemize}
    \item Los resultados, razonamientos (como en la pregunta 3) y manejo de unidades del Sistema Internacional fueron aplicados impecablemente en casi todos los ejercicios.
    \item Se recomienda \textbf{no omitir nunca los pasos de sustitución y desarrollo matemático}. En ciencias e investigación, es tan importante el resultado final como evidenciar claramente el método y la fórmula que permitieron llegar a él (como ocurrió en la Pregunta 5).
    \item Se incentiva a mantener este fantástico nivel de organización en sus futuros reportes y evaluaciones.
\end{itemize}

\end{document}